% section: 等差数列とその総和の導出

\section{等差数列とその総和の導出}\label{節：等差数列とその総和の導出}
  \noindent\textbf{【本編で使う場面】}\par
  等差数列の和を,並べ方の工夫から導く.\sectionref*{節：階差型}で一次式を足すときに使う.

  初項が$a$,公差が$d$である等差数列
  \[
    a,\,a+d,\,a+2d,\,\ldots,\,a+(n-1)d
  \]
  の初項から第$n$項までの総和を考える.
  その総和を$S_n$とし,第$n$項を$l$とすると,
  \[
    l=a+(n-1)d
  \]
  である.

  \subsection{総和の公式}\label{節：等差数列とその総和の導出：総和の公式}
    等差数列の初項から第$n$項までの総和は,
    \[
      S_n
      =
      \frac{n}{2}\{2a+(n-1)d\}
    \]
    である.
    第$n$項$l=a+(n-1)d$を用いれば,
    \[
      S_n=\frac{n}{2}(a+l)
    \]
    とも表せる.
    つまり,総和は項数と初項・末項の平均との積である.

  \subsection{公式の導出}\label{節：等差数列とその総和の導出：公式の導出}
  \BookFigArithmeticSum

    図は公差$d>0$の場合に,各項の増分
\BookMathLayout{$d,\,2d,\,\ldots,\,(n-1)d$と,それを逆順に並べたものを上下に重ねた模式図である.}{\\
$d,\,2d,\,\ldots$,$(n-1)d$と,それを逆順に並べたものを上下に重ねた模式図である.}
    上下を組み合わせた高さはどの列も$nd$になる.
    つまり,
    \[
      2\sum_{k=1}^{n-1}kd=n(n-1)d
    \]
    であり,共通の初項$a$を$n$回加えれば,
    $S_n=na+\displaystyle\frac{n(n-1)}{2}d$となる.
    以下では,同じ工夫を元の等差数列そのものに適用する.


    総和を順に並べた式と,逆順に並べた式を書く.
    \[
      \begin{array}{rcl}
        S_n
          &=&
          a+(a+d)+(a+2d)\\
          &&\qquad +\cdots+\{a+(n-1)d\}\\
        S_n
          &=&
          \{a+(n-1)d\}+\{a+(n-2)d\}\\
          &&\qquad +\cdots+(a+d)+a
      \end{array}
    \]
    これらを項ごとに足すと,各列の和は$2a+(n-1)d$となる.
    そのような列が$n$個あるので,
    \[
      2S_n=n\{2a+(n-1)d\}
    \]
    である.
    両辺を$2$で割れば,
    \[
      S_n=\frac{n}{2}\{2a+(n-1)d\}
    \]
    を得る.

    この導出では,同じ数列を逆順に並べて加えることで,各列の和を一定にしている.

